% =====================================================================
% Appendix E --- Closures and the neo-Keynesian mapping
% ---------------------------------------------------------------------
% Ported from paper/to_evaluate/closures_and_demand_shocks.md (2026-09-18).
% This appendix is NEW relative to the outline's A-D plan; it carries the
% CGE-vs-IO cleavage, the FIDELIO mapping and the closure taxonomy that
% Section 3 introduces.
%
% *** REVIEWER FLAG (Stage-3 must-not-appear list) ***
% The closing subsection's dominance wording ("the wage regime is the
% first-order determinant of the aggregate multiplier"; "intersectoral
% mobility eta produces second-order reallocation effects") sits close to
% the retired eta-dominance claims. Check it against the must-not-appear
% list before this text enters the manuscript. The descriptive mapping
% (Sections 1-4) is uncontroversial; the summary's two claims are not.
% =====================================================================
\section{Closures and the neo-Keynesian mapping}\label{app:mapping}

\noindent This appendix maps the CGE/IO cleavage and the neo-Keynesian
alternative --- the EU JRC FIDELIO model --- onto the closure matrix. It
supports Section~\ref{sec:two-questions}.

\subsection{The core theoretical cleavage: demand shocks in CGE versus IO models}\label{app:cleavage}

A central question in applied input--output modeling --- especially for
evaluating green public investment programmes such as the EU Recovery and
Resilience Facility or large-scale building renovation --- is whether an
exogenous demand stimulus:

\begin{enumerate}
\item \textbf{expands real aggregate output} by mobilising unutilised
  resources (the Keynesian / Leontief perspective), or
\item \textbf{merely reallocates scarce resources}, causing factor prices to
  rise and crowding out private activity with virtually zero net change in
  aggregate real GDP (the canonical neoclassical CGE perspective).
\end{enumerate}

\paragraph{Why standard CGE models produce minimal aggregate effects from demand shocks.}
In textbook Walrasian CGE models: (i) \emph{full factor utilisation} ---
aggregate primary factor endowments $(\bar L, \bar K)$ are exogenously given
and fully employed; (ii) \emph{instantaneous market clearing} --- prices and
wages are perfectly flexible and clear all commodity and factor markets
instantaneously; (iii) \emph{Say's Law operates} --- output is
supply-determined by the production possibility frontier
$Y = F(K, L, \mathbf{X})$; (iv) \emph{crowding out} --- an exogenous increase
in public investment raises demand in targeted sectors, but because total
labour and capital cannot expand, the extra factor demand bids up wages and
capital rentals, and the relative price increase squeezes the margins of
non-shocked sectors, crowding out private consumption, private investment or
net exports. The aggregate real GDP response is negligible
($\Delta Y \approx 0$).

\subsection{The macroeconometric alternative: the EU JRC FIDELIO model}\label{app:fidelio}

The \textbf{FIDELIO} model (\emph{Fully Interregional Dynamic Econometric
Long-term Input--Output}; JRC / DG GROW; Rocchi et al., 2019, 2025) was built
specifically to overcome this limitation. Rather than adopting a Walrasian
general-equilibrium core, FIDELIO is a \emph{dynamic macroeconometric,
multi-sectoral Neo-Keynesian} model grounded in the official Eurostat/JRC
FIGARO inter-country input--output tables (64 NACE industries, 46
countries/regions). It generates positive, non-trivial real output and
employment multipliers from demand shocks through five structural features:

\begin{enumerate}
\item \textbf{Involuntary unemployment and the wage-setting curve.} Wages do
  not clear the labour market. The baseline economy features involuntary
  unemployment ($u > 0$). Nominal wages adjust sluggishly along an
  econometrically estimated \emph{wage curve} (Blanchflower--Oswald
  tradition):
  \[
  \Delta \ln w_{i,r,t} = f(\text{inflation}_{t-1},\ \text{productivity}_{t},\ u_{r,t-1}).
  \]
  Because labour supply is not vertical at baseline, firms can hire idle
  workers to expand production without hitting an immediate supply wall.
\item \textbf{Mark-up pricing and sluggish cost pass-through.} Product prices
  are set via markups over unit costs (derived from flexible Translog cost
  functions over capital, labour, energy and intermediate inputs) rather than
  jumping to Walrasian clearing levels. Output in the short-to-medium run is
  \emph{demand-determined}.
\item \textbf{Keynesian multiplier loops (the ECM consumption block.)}
  Household consumption follows a dynamic error-correction model tied to
  current real disposable income rather than to an intertemporal Euler
  equation:
  \[
  \Delta G \Rightarrow \uparrow Y \Rightarrow \uparrow L
  \Rightarrow \uparrow (w\cdot L) \Rightarrow \uparrow C^{\text{household}},
  \]
  inducing a secondary Keynesian expenditure loop (a Type-II multiplier),
  yielding aggregate output multipliers between 1.2 and 1.8 in JRC policy
  simulations.
\item \textbf{Dynamic capital accumulation (endogenous capacity).} Capital is
  not a fixed static endowment; investment responds dynamically through an
  accelerator / user-cost formulation, so sustained demand increases capacity
  utilisation and expected returns, prompting firms to invest and expand
  capital stocks over time.
\item \textbf{Medium-to-long-run re-equilibration.} As unemployment drops and
  capacity tightens, the wage curve drives up wages and unit costs; the
  resulting rise in output prices erodes export competitiveness, generating
  gradual crowding out toward a medium-run NAIRU --- unless the stimulus was
  in public capital or R\&D, in which case it permanently shifts potential
  output outward.
\end{enumerate}

\subsection{Mapping into the \texttt{BeyondHulten} $5\times3$ closure matrix}\label{app:matrix-map}

The platform organises macroeconomic closures across two dimensions (five
labour closures $\times$ three financing closures).

\paragraph{The five labour closures.}
\begin{enumerate}
\item \textbf{ALPHA} (full-employment mobile CGE benchmark):
  $\sum_i L_i = \bar L$, single flexible wage $w$, cost-minimising
  intersectoral allocation ($\eta = 1$).
\item \textbf{BF} (immobile-labour endpoint selector): $\eta = 0$ freezes
  sectoral labour quantities to baseline shares $L_i = \bar L_i$; with the
  sectoral-wage endpoint (ADR-0020) the wage vector is free and the frozen
  allocation carries a price channel.
\item \textbf{BETA} (elastic labour supply on the real wage):
  $\sum_i L_i = \bar L\,[(w/P)/(w_0/P_0)]^{\eta_s}$, deflated by the CPI.
\item \textbf{GAMMA} (fixed real wage, uncapped extensive margin):
  $w/P = \bar w = 1$, employment $L = \sum_i L_i$ completely
  demand-determined.
\item \textbf{DELTA} (Leontief IO corner): GAMMA ($w/P = 1$) combined with
  the zero-substitution limit of the CES production core
  ($\theta, \epsilon, \sigma \to 0^+$), reproducing the fixed-price Leontief
  multiplier.
\end{enumerate}
(Complementarity extension \textbf{ZETA}: $0 \le \bar L - L \perp w/P -
\bar\omega \ge 0$, regime-switching between GAMMA slack and ALPHA full
employment; recorded as an idea, not a closure.)

\paragraph{The three financing closures.}
\begin{enumerate}
\item \textbf{F1} (preference reallocation / demand tilt): budget-neutral
  composition shift within fixed household expenditure
  ($\sum_i p_i c_i = E_h$); no net fiscal injection.
\item \textbf{F2} (tax-financed public investment): balanced government
  budget ($\sum_i p_i g_i = T(p)$); autonomous public investment financed by
  an explicit domestic tax counterparty.
\item \textbf{F3} (external / debt-financed programme): autonomous public
  spending financed externally or via debt ($\sum_i p_i g_i = F$), with
  imports and external balances absorbing the deficit.
\end{enumerate}

\paragraph{Matrix location of models.}
\begin{table}[htbp]
\centering
\small
\begin{tabular}{l l l l}
\toprule
Labour row & F1 (demand tilt) & F2 (tax-financed) & F3 (debt / external) \\
\midrule
ALPHA (full $L$)   & sectoral reallocation & full crowding out & crowding out ($w$ bids up) \\
BF (immobile)      & bottlenecks / wedges  & contractionary wedge & bottleneck limits \\
BETA (elastic)     & moderate shift        & mild expansion & expansion via $w/P$ \\
GAMMA (fixed $w$)  & demand-driven shift   & balanced-budget multiplier & Keynesian multiplier \\
DELTA (IO/Leontief)& IO composition shift  & Leontief balanced multiplier & pure IO multiplier \\
\bottomrule
\end{tabular}
\caption{Where the standard CGE and the FIDELIO model sit in the matrix.}
\label{tab:matrix-map}
\end{table}

\begin{itemize}
\item \textbf{Where standard CGE sits:} ALPHA-F1 or ALPHA-F2 --- fixed
  aggregate labour $\bar L$ and flexible wages force crowding out; real output
  effects are close to zero.
\item \textbf{Where FIDELIO sits:} in its static / basic IO mode (modules
  1--2) at DELTA-F3 (a pure Leontief input--output multiplier driven by
  external funding); in its short-run dynamic mode at GAMMA-F3 (fixed real
  wage, unconstrained extensive margin, autonomous injection), activating the
  full Keynesian respending multiplier; and, in the medium-to-long run, moving
  dynamically from GAMMA-F3 toward BETA-F3 as the wage curve drives real wage
  growth.
\end{itemize}

\subsection{How neo-Keynesian assumptions map into our closures}\label{app:neokeynesian}

In this platform, neo-Keynesian assumptions are not auxiliary parameters ---
they define the mathematical structure of the closure.

\paragraph{A. Wage rigidity (extensive margin versus wage inflation).}
In accordance with ADR-0014 and DE-0004, wage rigidity here is defined on the
\emph{real wage} ($w/P = \bar w$), deflated by the CPI numeraire ($P = 1$).
In GAMMA, fixing $w/P = 1$ removes the wage as an equilibrating price: the
labour-market equation changes from clearing $L = \bar L$ to an unconstrained
demand equation $L = \sum_i L_i(y_i, \mathbf{p})$, and the extensive margin of
employment absorbs the demand shock. Asymmetric rigidity (ZETA) follows the
spirit of Baqaee--Farhi (2022): downward wage rigidity creates a regime
switch --- demand deficits cause unemployment (GAMMA regime), while demand
surges that exhaust slack trigger wage inflation (ALPHA regime).

\paragraph{B. Price rigidity and fixed markups.} In standard New-Keynesian
models, Calvo or Rotemberg pricing prevents prices from adjusting
instantaneously to demand shocks. Here, under GAMMA, wages are pegged
($w = 1$) and the shock is demand-driven, so marginal costs remain roughly
constant and product prices exhibit high inertia ($\mathbf{p} \approx 1$). In
DELTA, the zero-substitution limit ($\theta, \epsilon, \sigma \to 0^+$)
prevents firms from adjusting technical factor proportions, so cost
pass-through becomes purely linear and additive, reproducing the classic
fixed-price Leontief/Robinson (2006) multiplier.

\paragraph{C. Non-Ricardian demand and disposable-income loops.} Standard CGE
models enforce intertemporal optimisation in which households anticipate
future tax liabilities, neutralising debt-financed spending. Here the
consumption block routes household income $E_h = (1-\tau) w \sum_i L_i +
\text{capital income}$; under F3 the absence of an immediate tax hike allows
the newly generated wage income $w\cdot\Delta L$ to fuel private consumption
directly, generating an endogenous Keynesian multiplier.

\subsection{Summary and methodological takeaway}\label{app:summary}
% *** Flagged above: verify against the must-not-appear list. ***

\begin{enumerate}
\item \textbf{The binary wage regime dominates.} The literature often debates
  the elasticity of factor substitution or intersectoral mobility ($\eta$).
  As shown in \texttt{definitive\_guide.md} and our runs, intersectoral
  mobility produces second-order reallocation effects, whereas the wage regime
  (flexible ALPHA versus sticky GAMMA) is the first-order determinant of the
  aggregate multiplier.
\item \textbf{Reconciling applied discrepancies.} When policy studies using
  macroeconometric models (FIDELIO, E3ME) report output multipliers around
  $1.5$ for green investments while standard CGE models report zero or
  negative multipliers, the divergence is not an empirical dispute over green
  technology. It is a direct mathematical consequence of evaluating the shock
  at GAMMA/DELTA-F3 (slack resources, wage stickiness, debt financing) versus
  ALPHA-F2 (full employment, market clearing, domestic tax financing).
\end{enumerate}